% Generated draft: requires mathematical and source review.
\documentclass[11pt]{article}
\usepackage[margin=25mm]{geometry}
\usepackage{fontspec}
\setmainfont{Latin Modern Roman}
\usepackage{amsmath,amssymb,mathtools,unicode-math}
\setmathfont{Latin Modern Math}
\usepackage{enumitem,needspace,xurl,hyperref}
\setlength{\parindent}{0pt}
\setlength{\parskip}{5pt}
\setlength{\emergencystretch}{3em}
\newcommand{\N}{\mathbb{N}}
\newcommand{\Z}{\mathbb{Z}}
\newcommand{\Q}{\mathbb{Q}}
\newcommand{\R}{\mathbb{R}}
\newcommand{\C}{\mathbb{C}}
\newcommand{\F}{\mathbb{F}}
\newcommand{\A}{\mathbb{A}}
\newcommand{\PP}{\mathbb{P}}
\newcommand{\E}{\mathbb{E}}
\newcommand{\eps}{\varepsilon}
\newcommand{\dd}{\,\mathrm{d}}
\begin{document}
\section*{UnsolvedMath ID: 11000008}
\textbf{Source problem code:} catalogue entry
\subsection{Definitions and mathematical statement}
Let \(G=\operatorname{Mod}_g\) be the mapping class group in the source's genus and surface convention. In the usual closed oriented convention it is the group of orientation-preserving self-homeomorphisms of a genus-\(g\) surface modulo isotopy. For a finite generating set \(S\) of \(G\) and \(H\leq G\), form the Schreier graph on cosets \(H\backslash G\) with edges \(Hg\leftrightarrow Hgs\), \(s\in S\cup S^{-1}\). In the ordinary geometric relative-end convention, \(\operatorname{ends}(G,H)\) is the number of ends of this locally finite graph: ends are consistent choices of infinite components outside increasing finite vertex sets. The finite graph has zero ends. This number is independent of the finite generating set. The excerpt does not explicitly distinguish geometric relative ends from other algebraic relative-end invariants.

\emph{Question 2.1 (ends spectrum).} For finitely generated subgroups \(H<\operatorname{Mod}_g\), determine all possible values of
\[
 \operatorname{ends}(\operatorname{Mod}_g,H).
\]
The genus range, whether the subgroup notation is strict, and the relative-end convention must be fixed. The summary invokes the familiar \(0,1,2,\infty\) alternatives for ends of finitely generated groups, but that theorem does not by itself impose the same spectrum on general Schreier graphs of subgroup pairs. No such restriction is silently added.

The source states that the moduli space \(\mathcal M_g\) has one end, with connectedness of the curve complex as the key argument, and that the proof gives one-endedness of every cover; it refers to [FMa]. It expressly says that this does not directly explain the answer to the relative-end question. The source's geometric covering-space observation is therefore not identified with the Schreier-graph invariant.

The extraction then begins a separate discussion of commensurators. It calls subgroups \(\Gamma_1,\Gamma_2\leq A\) commensurable up to conjugacy if some \(a\in A\) makes \(a\Gamma_1a^{-1}\cap\Gamma_2\) finite index in both groups. It defines
\[
 \operatorname{Comm}_A(\Gamma)
 =\{a\in A:[\Gamma:\Gamma\cap a\Gamma a^{-1}]<\infty,
 [a\Gamma a^{-1}:\Gamma\cap a\Gamma a^{-1}]<\infty\}.
\]
It describes this as ignoring finite-index information beyond conjugacy and records Margulis's criterion that, for an irreducible lattice \(\Gamma\) in a semisimple Lie group \(A\),
\[
 [\operatorname{Comm}_A(\Gamma):\Gamma]=\infty
 \quad\Longleftrightarrow\quad\Gamma\text{ is arithmetic}.
\]
Its footnote specifies that ``semisimple'' means linear semisimple with no compact factors. These are contextual definitions and an attributed theorem with its source scope, not further ends-spectrum requirements or an independent verification of the lattice statement.
\subsection{Short English statement}
Determine the possible relative end numbers for finitely generated subgroups of the mapping class group. The source distinguishes this from the one-endedness of moduli-space covers; the standard end theorem for groups alone does not determine the Schreier-graph spectrum.
\subsection{Sources}
\begin{itemize}
\item Local UnsolvedMath registry, permanent ID 11000008, source code catalogue entry, Question 2.1, ``Ends spectrum,'' PDF page 22. No individual proposer field is supplied; the text and summary identify B. Farb's discussion.
\item Farb, \emph{Problems on Mapping Class Groups and Related Topics}. Supplied URL: \url{https://www.math.uchicago.edu/~farb/papers/mcgbook.pdf}; recorded access 2026-07-29; extraction pdf-font-structure. The edited volume contains 329 explicitly typeset problem, conjecture, and question headings across chapters, with repeated numbers disambiguated by page.
\item The statement cites [FMa] for moduli-space covers and attributes its trailing commensurator criterion to Margulis; full entries are absent. The summary refers to Stallings--Swan theory, Farb--Mosher, \emph{Geom. Topol.} (2002), and Durham--Taylor on stable subgroups, without complete citations. It describes the latter subgroup theories as related but not resolving the full spectrum.
\item Supplied status open and OPEN-TRIAGE. The summary reports no later resolution located and incorrectly suggests that the group-end alternatives alone specify the candidate pair spectrum; this implication is flagged above. Current primary-literature status and redistribution terms for the publicly accessible author/publisher copy remain NEEDS\_REVIEW.
\end{itemize}
\end{document}
